\section{Inherited operator drift and the full coding consequence}

This section reproduces the actual Cycle5 proof \cite{Cycle5}. Its hypotheses are verified by (28); no fresh-channel replacement is made.

\begin{lemma}[Homogeneous increments]\label{lem:grade}
For any charge-preserving $k$-call density $\rho$, its normal bath channel
has Kraus operators of homogeneous depth degree
$\Delta\in\{-k,\ldots,k\}$. Consequently there are positive effects
$A_\Delta$ such that
\[
 \sum_\Delta A_\Delta=I,\quad [A_\Delta,L]=0,\quad
 D_s(\rho)=\sum_\Delta\Delta A_\Delta,
 \tag{C1}
\]
and, as a positive quadratic-form identity,
\[
 \Phi_\rho^*e^{tL}
 =\left(\sum_\Delta e^{t\Delta}A_\Delta\right)e^{tL}.
 \tag{C2}
\]
\end{lemma}
\begin{proof}
Diagonalize $\rho$ in joint eigenvectors of $Q_k$, and take computational
output words $o$. Charge conservation gives each corresponding bath Kraus
operator degree $\Delta=q-N(o)$:
$[L,K]=\Delta K$. The exchange $R$ has degree zero and satisfies this
identity at the vacuum as well. Set $A_\Delta=\sum_{\deg K=\Delta}K^*K$.
Trace preservation proves their sum is $I$; homogeneity proves they commute
with $L$ and yields (C1)--(C2) on finite-coordinate vectors. Positive
quadratic forms and monotone limits extend the exponential identity.
\end{proof}

\begin{theorem}[Uniform operator drift gives actual reliable codes]
\label{thm:operatorcode}
Fix $k$, $0<\beta<1$, and a charge-preserving density $\rho$.
Suppose its exact drift obeys
\[
 P_{\ge k}D_s(\rho)P_{\ge k}\le-\kappa P_{\ge k},
 \qquad0<\kappa\le k.
 \tag{C3}
\]
Then one deterministic unassisted code sequence, independent of the bath
initializer, has rate at least $S(\rho)/k$ and pointwise vanishing
full-diamond error for every fixed normal initializer.
For the actual $\tau_\pi$, the product-block coherent-information limit
is also $S(\rho)/k$. No scalar depth chain or radial state preservation
is required.
\end{theorem}
\begin{proof}
Set
\[
 t=\kappa/(2k^2),\quad
 \lambda=1-\kappa^2/(4k^2),\quad
 C=\frac{e^{2kt}}{1-\lambda}.
 \tag{C4}
\]
Since $kt\le1/2$, for $|\Delta|\le k$ Taylor's inequality gives
$e^{t\Delta}\le1+t\Delta+k^2t^2$.
Multiply by each positive effect in (C1) and sum. On the high-depth sector
this is at most $\lambda I$. The low-depth sector reaches depth at most
$2k-1$, so (C2) gives the global form inequality
\[
 \Phi_\rho^*e^{tL}
 \le\lambda e^{tL}+e^{2kt}P_{<k}.
 \tag{C5}
\]
There are no low/high cross terms. Iteration yields
\[
 M_m\le\lambda^m M_0+C(1-\lambda^m).
 \tag{C6}
\]
For any coherent initializer supported in $P_{\le A}$,
\[
 \Pr(L_m>H)\le(1+C)e^{-t(H-A)}.
 \tag{C7}
\]

Let $d=\rank\rho\le3^k$, using only positive eigenvalues. A
maximum-probability exact spectral type $T_m$ of $\rho^{\otimes m}$
has probability $w_m\ge(m+1)^{-d}$, rank $D_m$, and
\[
 T_m/D_m\le\rho^{\otimes m}/w_m,\qquad
 \log_2D_m/m\longrightarrow S(\rho).
 \tag{C8}
\]
For completeness, with type counts $n_i$, maximality against a transfer
from count $i$ to $j$ gives
$n_i/p_i\le(n_j+1)/p_j$. Summing against $p_i$ shows
$n_j\ge mp_j-1$, hence every frequency tends to $p_j$.
The type scalar eigenvalue and the polynomial mass give the rank limit.
This is an operator statement in the actual eigenbasis, even if eigenvectors
are entangled across calls.

Choose
\[
 A_m=\lfloor\log_2(m+1)\rfloor,\quad
 H_m=A_m+\left\lceil
 \frac{(d+16)\ln(m+1)+\ln(1+C)}t\right\rceil,\quad
 \epsilon_m=(m+1)^{-16}.
 \tag{C9}
\]
Apply (C7) to the uniform state on $P_{\le A_m}$, and then (C8).
Mean overflow for the type subspace and that initial bath space is at
most $\epsilon_m$. The inherited joint compression lemma, reproduced
in Appendix~\ref{app:compression}, gives
\[
 K_m\ge\left\lceil
 \frac{\lfloor D_m/2\rfloor}{(r_{A_m}r_{H_m})^4}
 \right\rceil,\quad
 e_m=r_{A_m}\sqrt{2\epsilon_m}+(1-2\epsilon_m)^{-1/2}-1\to0.
 \tag{C10}
\]
It retains spectral expurgation, the $r^2$ positive-block factor,
the auxiliary CPTP normalization and its polar repair.
The actual encoder is isometric; no auxiliary normalization is implemented.
For every fixed normal $\tau$, the full-diamond error is at most
\[
 \min\{2,\ 4e_m+2\sqrt{1-\Tr(P_{\le A_m}\tau)}\}\longrightarrow0.
 \tag{C11}
\]
In half-diamond convention this is
$\min\{1,2e_m+\sqrt{1-\Tr(P_{\le A_m}\tau)}\}$.
This follows by comparing $\tau$ outside the type domination with its
normalized coordinate truncation. Thus no type-probability factor multiplies
an uncontrolled initial tail. Since $\log r_{A_m}r_{H_m}=O(\log m)$
for this fixed device and block, (C10) proves the rate.

For $\tau_\pi$, $t\le1/2<\ln3$ ensures a finite initial exponential
moment. Equations (C6)--(C7) bound mean depth uniformly. The bath Gibbs
density $g=4^{-L}/4$ has trace one and
$-\log_2g=2L+2$, so $S(E_m)\le2\E L_m+2$.
This entropy-energy inequality also follows by finite cutoffs and monotone
limits; finite mean guarantees finite entropy. The initial bath entropy
is finite. Unitary entropy conservation, subadditivity and Araki--Lieb give
\[
 mS(\rho)-2S(E_m)\le\Ic(\rho^{\otimes m},\mathcal N_{km}^{\tau_\pi})
 \le mS(\rho).
 \tag{C12}
\]
Divide by $km$. Extra data-independent calls may fill incomplete final
blocks; all count and do not alter already emitted output marginals.
\end{proof}

The initializer-independent code quantifier is pointwise in the normal initializer.
It is distinct from the uniform all-normal upper bound (4).
No invariant initializer is prepared. The boundary term in (C5) is essential:
a global strict exponential contraction would contradict trace preservation
and $e^{tL}\ge I$.

